% Suggested changes for chapter 04_WindTurbineControl_02_notBold.tex (main 13a333a):
% numbers that change with the evaluation from 60 s instead of 66 s (IPC/AALC runs repeated with
% the FAST output from 60 s; tables figures/figIPCThesis/tab_IPC_summary.tex and
% tab_AALC_summary.tex and the bar/Pareto figures are updated). All other numbers in the text
% of the IPC, AALC and comparison sections remain the same at the precision given.

% ---- original line 854 (after: "")
% IPC table above rated: C_DEL_t 0.997, 0.999, 1.005 (Verif. 2 was 1.01 with 66 s)
% OLD: Above rated, the blade DEL decreases by 19\% to 22\%, and from 18\,m/s on by 23\% to 24\%, for all three datasets. The reduction grows with the wind speed: for the second verification set of turbulence class~B, it increases from 12\% to 15\% between 12 and 16\,m/s to 24\% to 27\% from 20\,m/s on. The gains optimized at a single wind speed hence carry over to the whole of Region~III and to both verification sets, so that no gain scheduling is required above rated. The tower base DEL is not changed considerably, with criteria between 0.99 and 1.01; the largest increase of a single wind speed is 3\% at 18\,m/s for turbulence class~B. The tower top DEL is reduced slightly, by 2\% to 4\% above rated and by 3\% to 6\% at 18\,m/s. Close to rated wind speed, the mean power decreases by up to 0.6\% and its standard deviation increases slightly, so that the power criteria above rated are 1.00 and 1.01 to 1.02. The pitch rate reaches the limit of 8\textdegree/s at all wind speeds, and the actuator power is increased by a factor of 5.4 to 5.7 above rated and of 5.8 to 6.0 from 18\,m/s on. Compared with the AALC (Table~\ref{tabone}), the IPC achieves about the same blade DEL reduction above rated (AALC: 19\% to 23\%), but at a clearly higher actuator power (AALC: factor 3.9 to 4.0); from 18\,m/s on, the AALC reduces the blade DEL slightly more (22\% to 26\%).
Above rated, the blade DEL decreases by 19\% to 22\%, and from 18\,m/s on by 23\% to 24\%, for all three datasets. The reduction grows with the wind speed: for the second verification set of turbulence class~B, it increases from 12\% to 15\% between 12 and 16\,m/s to 24\% to 27\% from 20\,m/s on. The gains optimized at a single wind speed hence carry over to the whole of Region~III and to both verification sets, so that no gain scheduling is required above rated. The tower base DEL is not changed considerably, with criteria of 1.00 for all three datasets; the largest increase of a single wind speed is 3\% at 18\,m/s for turbulence class~B. The tower top DEL is reduced slightly, by 2\% to 4\% above rated and by 3\% to 6\% at 18\,m/s. Close to rated wind speed, the mean power decreases by up to 0.6\% and its standard deviation increases slightly, so that the power criteria above rated are 1.00 and 1.01 to 1.02. The pitch rate reaches the limit of 8\textdegree/s at all wind speeds, and the actuator power is increased by a factor of 5.4 to 5.7 above rated and of 5.8 to 6.0 from 18\,m/s on. Compared with the AALC (Table~\ref{tabone}), the IPC achieves about the same blade DEL reduction above rated (AALC: 19\% to 23\%), but at a clearly higher actuator power (AALC: factor 3.9 to 4.0); from 18\,m/s on, the AALC reduces the blade DEL slightly more (22\% to 26\%).

% ---- original line 1044 (after: "")
% AALC table above rated: C_Act 3.86, 3.85, 4.02 (as already stated in the IPC paragraph: "AALC: factor 3.9 to 4.0")
% OLD: The evaluation of the power criteria confirms that the AALC leads only to negligible changes in both power mean and variance above rated. Below rated, the mean power is decreased by 2\% to 3\%, as the flap and pitch deflections reduce the aerodynamic efficiency, cf.\ Figure~\ref{fig:OLval}. The additional rotational movement of the blade demanded by IPC is the only criterion leading to an unavoidable increase, similar to the two controller schemes discussed in previous sections: above rated, the actuator power is increased by a factor of 3.8 to 4.0, and from 18\,m/s on by a factor of 3.9 to 4.2. Below rated, the baseline pitch angle is almost constant, so that the baseline actuator power is close to zero and no meaningful ratio is obtained. To assess the overall performance of AALC relative to the other discussed control strategies, the following section presents a comparison of all three investigated algorithms on load reduction and  power standard deviation at different pitch rate limits.
The evaluation of the power criteria confirms that the AALC leads only to negligible changes in both power mean and variance above rated. Below rated, the mean power is decreased by 2\% to 3\%, as the flap and pitch deflections reduce the aerodynamic efficiency, cf.\ Figure~\ref{fig:OLval}. The additional rotational movement of the blade demanded by IPC is the only criterion leading to an unavoidable increase, similar to the two controller schemes discussed in previous sections: above rated, the actuator power is increased by a factor of 3.9 to 4.0, and from 18\,m/s on by a factor of 3.9 to 4.2. Below rated, the baseline pitch angle is almost constant, so that the baseline actuator power is close to zero and no meaningful ratio is obtained. To assess the overall performance of AALC relative to the other discussed control strategies, the following section presents a comparison of all three investigated algorithms on load reduction and  power standard deviation at different pitch rate limits.

% ---- original line 1103 (after: "by 33\%, and requires the highest actuator power of all controllers wi")
% smooth design, 8 deg/s: C_dP 0.762 (66 s: 0.767)
% OLD: of the PI baseline. Filters were applied inside the qLMPC to make it less aggressive. The results are highlighted here in pink: the smooth design trades part of the power smoothing, a reduction by 23\%,
of the PI baseline. Filters were applied inside the qLMPC to make it less aggressive. The results are highlighted here in pink: the smooth design trades part of the power smoothing, a reduction by 24\%,

% ---- original line 1116 (after: "The middle panels show the IPC and the AALC, which act on the blade lo")
% IPC, 8 deg/s: C_DEL_b 0.798 (66 s: 0.790)
% OLD: 4\textdegree/s to 0.79 at 8\textdegree/s. A greater rate limit of 13\textdegree/s brings no further benefit,
4\textdegree/s to 0.80 at 8\textdegree/s. A greater rate limit of 13\textdegree/s brings no further benefit,

% ---- original line 1126 (after: "about the actuator power of the IPC and the AALC. It almost reaches th")
% IPC, 8 deg/s: C_DEL_b 0.798; smooth design + IPC: C_DEL_b 0.852, C_dP 0.751 (66 s: 0.755)
% OLD: IPC, 0.85 compared with 0.79, and in addition reduces the power standard deviation by about 24\%,
IPC, 0.85 compared with 0.80, and in addition reduces the power standard deviation by about 25\%,
